\documentclass[11pt]{article}

\usepackage{common/tonas-papers}

\graphicspath{{./images/}}

\begin{document}

\makeboxedtitle{No countable basis for Borel digraphs of dichromatic number at least three}{
    Tonatiuh Matos-Wiederhold
}{}{
    I prove that the Borel directed graphs admitting a partition into two Borel acyclic sets of vertices form a $\mathbf\Sigma^1_2$-complete set; equivalently, that the set of Borel digraphs with Borel dichromatic number at least~$3$ is $\mathbf\Pi^1_2$-complete.
    It follows that no countable family of Borel digraphs can serve as a basis for this class under Borel homomorphism and, more generally, that any basis must be at least as complex as~$\mathbf\Pi^1_2$.
    The proof lifts the classical NP-completeness reduction of Bokal, Fijavž, Juvan, Kayll, and Mohar to the Borel setting, using the coding framework of Thornton.
    This in contrast with the recent result of Raghavan and Xiao, who constructed a continuum-size basis of Borel directed graphs for Borel directed graphs of \emph{uncountable} Borel dichromatic number, which adds to the picture initiated by Todorčević and Vidnyánszky for the undirected case.
}

\section{Introduction}\label{sec:borel-dichrom-intro}



\end{document}